\section{Damping the kernel}\label{cg:sec:kernel}

The asymmetry of Section~\ref{cg:sec:cmfamily} might be blamed on working with $g$ rather than the symmetric theta representation. This section applies the same damping, with Gauss weights, to the coefficients in Riemann's kernel. The mechanism is that of Chapter~\ref{chap:miyatake} (Lemma~\ref{mi:lem:kink}, Theorem~\ref{mi:thm:kink}, Propositions~\ref{mi:prop:kinkbudget} and~\ref{mi:prop:reweight}, with friable weights and the same proof). Write
\[
\Phi_n(u)=\bigl(2\pi^2n^4e^{9u/2}-3\pi n^2e^{5u/2}\bigr)\exp\bigl(-\pi n^2e^{2u}\bigr),\qquad \Phi=2\sum_{n\ge1}\Phi_n,
\]
so that $\Xi(t)=2\int_0^\infty\Phi(u)\cos(ut)\,du$, $\Phi>0$ and $\Phi$ is even; the evenness is the theta transformation formula. For $\delta\ge0$ let $\Phi_\delta=2\sum_ne^{-\pi n^2\delta}\Phi_n$ on $[0,\infty)$ and $\Xi_\delta(t)=2\int_0^\infty\Phi_\delta(u)\cos(ut)\,du$. With $x=e^{2u}$ this is the incomplete gamma form of $\Xi$ (Proposition~\ref{mi:prop:incgamma}) with weights $e^{-\pi n^2\delta}$, less the constant $\sum_n(1-e^{-\pi n^2\delta})(4\pi n^2-1)e^{-\pi n^2}$ (two integrations by parts, as in Section~\ref{mi:sec:approx}), the analogue of the factors $(1-\epsilon)^n$ in the paper of Comtet and Majumdar (their B6, not checked). Since $\Phi_\delta$ decays like $\exp(-\pi e^{2u})$, each $\Xi_\delta$ is even, real on $\R$ and entire of order one, $\Xi_\delta\to\Xi$ locally uniformly by dominated convergence, and $\Xi_\delta(0)>0$ because $\Phi_n>0$ on $[0,\infty)$; so $\Xi_\delta$ satisfies the hypotheses of Lemma~\ref{cg:lem:lpggc}. For complex $t$ the transform is $\Xi_\delta(t)=\sum_ne^{-\pi n^2\delta}[B_n(t)+B_n(-t)]$, where $B_n(t)=2\pi^2n^4\Gamma(z_9,\pi n^2)(\pi n^2)^{-z_9}-3\pi n^2\Gamma(z_5,\pi n^2)(\pi n^2)^{-z_5}$ with $z_9=\frac94+\frac{it}2$, $z_5=\frac54+\frac{it}2$ and $\Gamma(\cdot,\cdot)$ the upper incomplete gamma function; on the real axis $B_n(-t)=\overline{B_n(t)}$ and the sum is $2\operatorname{Re}B_n(t)$, the form usually printed, which is not the analytic continuation off the axis.

\begin{theorem}\label{cg:thm:kink}
For every $\delta>0$, $\Xi_\delta$ has only finitely many real zeros and infinitely many nonreal zeros. Consequently $\Xi_\delta(0)/\Xi_\delta(i\sqrt w)$ is not the Laplace transform of a GGC.
\end{theorem}

\begin{proof}
A direct computation gives $\Phi_n'(0)=c_n:=\pi n^2e^{-\pi n^2}\bigl(15\pi n^2-\tfrac{15}2-4\pi^2n^4\bigr)$. For $n=1$ the bracket is $15\pi-7.5-4\pi^2=0.1455>0$ (computed; the source printed $0.146$); for $n\ge2$ it equals $-n^2(4\pi^2n^2-15\pi)-7.5<0$ because $4\pi^2n^2>15\pi$. Evenness of $\Phi$ gives $\sum_nc_n=\Phi'(0)/2=0$. Hence for decreasing weights $w_n=e^{-\pi n^2\delta}$,
\[
\tfrac12\Phi_\delta'(0)=\sum_nw_nc_n=w_1c_1+\sum_{n\ge2}w_nc_n>w_1\sum_{n\ge1}c_n=0,
\]
since $w_n<w_1$ and $c_n<0$ for $n\ge2$. Two integrations by parts and the Riemann Lebesgue lemma, as in Lemma~\ref{ev:lem:kink} and Theorem~\ref{mi:thm:kink}, give that $t^2\Xi_\delta(t)\to-2\Phi_\delta'(0)<0$, so $\Xi_\delta<0$ for all large real $t$ and has finitely many real zeros. If $\Xi_\delta$ had finitely many zeros in all, Hadamard's theorem would give $\Xi_\delta=Pe^{at}$ with $P$ a polynomial; evenness and reality force $a=0$, and a nonzero polynomial cannot tend to zero along $\R$. So there are infinitely many zeros, all but finitely many nonreal, and Lemma~\ref{cg:lem:lpggc} applies.
\end{proof}

Damping the $n$ sum destroys the evenness of $\Phi$ at $u=0$, so the even extension $\Phi_\delta(|u|)$ has a kink, which produces a nonoscillating $t^{-2}$ tail that swamps the exponentially small oscillation $|\Xi(t)|\approx t^{7/4}e^{-\pi t/4}$ that carries the zeros. Heuristically the real zeros stop near the height $T(\delta)$ where these balance, $\frac\pi4T-\frac{15}4\ln T\approx\ln(1/\delta)$; for $\delta=10^{-3}$ this gives $T\approx24$ (heuristic).

\begin{corollary}[Modular constraints]\label{cg:cor:modular}
Let $\Phi_w=2\sum_nw_n\Phi_n$ with real weights such that $\sum_n|w_n|n^Ae^{-\pi n^2}<\infty$ for every $A\ge0$, and $\Xi_w(t)=2\int_0^\infty\Phi_w(u)\cos(ut)\,du$. If $\Xi_w\in\CGLP$ and $\Xi_w$ is not of the form $e^{-at^2}$ times a polynomial, then $\sum_nw_n\Phi_n^{(2j+1)}(0)=0$ for every $j\ge0$.
\end{corollary}

\begin{proof}
If the first nonzero odd derivative is of order $2j+1$, repeated integration by parts gives $\Xi_w(t)\sim2(-1)^{j+1}\Phi_w^{(2j+1)}(0)\,t^{-2j-2}$, of eventually constant sign, and the argument of Theorem~\ref{cg:thm:kink} applies (an element of $\CGLP$ with finitely many zeros is $e^{-at^2+bt}P$).
\end{proof}

For $w_n\equiv1$ these linear constraints are the Taylor coefficients at $x=1$ of Jacobi's identity $\theta(1/x)=\sqrt x\,\theta(x)$. The proof of Theorem~\ref{cg:thm:kink} shows that the first constraint fails for every weight sequence with $w_1\ge w_n\ge0$ for $n\ge2$ and at least one strict inequality (under the summability condition of Corollary~\ref{cg:cor:modular}): subtracting $w_1\sum c_n=0$ gives $\sum_{n\ge2}(w_n-w_1)c_n>0$. This covers the usual Abel, Gauss and Lerch weights; it is not a statement about arbitrary deformations of the coefficients.

\subsection{Symmetrising does not help}

One might restore evenness by defining $\Phi_\delta$ on the whole line through the $n$ sum and taking the Fourier transform over $\R$. Since $\int_\R\Phi_n(u)e^{iut}du=\frac14s(s-1)\pi^{-s/2}\Gamma(s/2)\,n^{-s}$ with $s=\frac12+it$, this gives
\[
\Xi^\natural_\delta(t)=\tfrac12\bigl[\xi^G_\delta(\tfrac12+it)+\xi^G_\delta(\tfrac12-it)\bigr],\qquad \xi^G_\delta(s)=\tfrac12s(s-1)\pi^{-s/2}\Gamma(s/2)\sum_{n\ge1}n^{-s}e^{-\pi n^2\delta},
\]
the real part of the Gauss damped completed zeta function on the real axis and its continuation elsewhere. Now $\Phi_\delta(u)\sim-6\pi\bigl(\sum_n n^2e^{-\pi n^2\delta}\bigr)e^{5u/2}$ as $u\to-\infty$, only exponential decay, because the double exponential decay of $\Phi$ at $-\infty$ is again the theta formula. Correspondingly $\xi^G_\delta$ has a pole at $s=-2$ and $\Xi^\natural_\delta$ has poles at $t=\pm\frac52i$. The reciprocal $1/\Xi^\natural_\delta(i\sqrt w)$ then vanishes at $w=25/4$, exactly as in Proposition~\ref{cg:prop:obstructions}(c).

\subsection{The numerics of the source}

The source tabulates zeros of the damped kernels, from mpmath at 18 to 30 digits with sign changes on a grid and an argument principle count on the box $0<\operatorname{Re}t<T$, $|\operatorname{Im}t|<4$. Apart from the three values of $\Phi_\delta'(0)$ below, none of it was recomputed and all of it is unchecked. For $\delta=0.05$ and $\delta=0.01$ the source finds one real zero in $(0,60)$ (at $14.1002$ and $14.1263$) and three zeros in the box, against $13$ for $\Xi$ itself; for $\delta=10^{-3}$ a scan to height $110$ finds three real zeros ($14.1338$, $21.0755$, $24.5364$); the ratio $t^2\Xi_\delta(t)/\{-2\Phi_\delta'(0)\}$ at $t=100$ is $0.973$ for three values of $\delta$; and $\frac12\Phi_\delta'(0)=1.85\times10^{-4}$, $1.72\times10^{-3}$, $6.34\times10^{-3}$ at $\delta=10^{-3},10^{-2},5\times10^{-2}$, the last three agreeing with a hand recomputation from $c_1=0.019749$, $c_2=-0.019749$, $c_3=-4.1\times10^{-8}$. For the symmetric damping $e^{-\delta\cosh2u}$ (the kernel $e^{-2\pi\cosh2u}$ of P\'olya's Bessel function is treated in Section~\ref{ch:biane} and Proposition~\ref{prop:biane:damping}, and the cosh ladder in Chapter~\ref{chap:miyatake}) the source finds that all zeros below height $45$ stay real for $\delta=0.5$ and $2$, eight and six of them, moving up and thinning out as $\delta$ grows (unchecked). If true this separates damping that respects Jacobi's identity from damping that does not, but it concerns low zeros only.

\section{The de Bruijn Newman family and what the identity contributes}\label{cg:sec:dbn}

By Corollary~\ref{cg:cor:modular} an admissible deformation must keep the even extension of $\Phi$ smooth at $u=0$; multiplying by a smooth even function of $u$ whose derivatives grow at most like $e^{Cu^2}$ does this and keeps the double exponential decay. The classical example is the de Bruijn Newman flow of Section~\ref{ch:dbn},
\[
\begin{gathered}H_t(z)=\int_0^\infty e^{tu^2}\,\Phi_{\mathrm{dBN}}(u)\cos(zu)\,du,\\ \Phi_{\mathrm{dBN}}(u)=\tfrac12\Phi(2u),\\ H_0(z)=\tfrac18\Xi(z/2).\end{gathered}
\]
The thresholds for $H_t$ are those of Section~\ref{ch:dbn}. The source writes the flow with $e^{\lambda u^2}\Phi(u)$ in the variable $z$, where $H_0=\Xi$; in that normalisation $H_\lambda(z)=8H_{4\lambda}(2z)$, the time scales by $4$, and the literature thresholds become $\lambda\ge\frac18$, $0.055$ and $0.05$. The source's statement with $\lambda\ge0.2$ is therefore true (zeros stay real forward in time) but is not the literature result it cites, and the Polymath bound is $0.22$, not $0.2$. The theorem is Theorem~\ref{rz:thm:dbn}: for $t\ge\Lambda$ the reciprocal $R_t=H_t(0)/H_t(2i\sqrt s)$ is a GGC transform with Thorin measure $\sum_k\delta_{(z_k(t)/2)^2}$ where $z_k(t)>0$ are the zeros of $H_t$, unconditionally for $t\ge0.2$, and RH holds if and only if $R_t$ is a GGC for arbitrarily small $t>0$, that is by Rodgers and Tao if and only if $\Lambda=0$. The source's statement that the transforms converge as $\lambda\downarrow0$ if and only if RH holds is false in the direction convergence implies RH: $H_t(i\sqrt w)>0$ for every $t$, and dominated convergence gives $H_t(0)/H_t(i\sqrt w)\to H_0(0)/H_0(i\sqrt w)$ for every $w>0$ without any hypothesis.

The Gaussian multiplier is $\exp(-t\,d^2/dz^2)$ acting on $H_0$, which preserves real zeros forward in time; a decreasing multiplier is the analogue of $t<0$, and Rodgers and Tao prove that every $t<0$ produces nonreal zeros. For the Gaussian family with $t<0$ and for the coefficient damping of Theorem~\ref{cg:thm:kink} this is a theorem. That damping is the wrong direction for every decreasing multiplier, whether applied to $n$ or to $u$, is a heuristic and is not proved here (the cosh numerics above, if confirmed, show low zeros staying real), and nothing established in this chapter decides the Laguerre P\'olya property of other decreasing multipliers.

What the identity of Comtet and Majumdar contributes is a probabilistic reading of the signed object that must become positive. By Proposition~\ref{cg:prop:signed}, $\zeta$ is the Mellin transform of the negative of the signed $k$ measure $dt-d\lfloor t\rfloor$ of the uniform law (equation \eqref{cg:eq:zetamellin}). Under RH, $\Xi(0)/\Xi(i\sqrt w)$ is the GGC transform with Thorin measure the counting measure of the squared ordinates. The identity therefore links two counting measures: integers minus Lebesgue on the arithmetic side, zeros on the spectral side. For real $s\in(0,1)$, $g<0$ and $\sin(\pi s/2)>0$ give $\zeta(s)<0$ by \eqref{cg:eq:mellin}, which is the statement that $\zeta$ has no real zeros there.

\begin{problem}\label{cg:prob:explicit}
Find the analogue, for the explicit formula, of the convexity and Poisson summation argument of Theorem~\ref{cg:thm:uniform}: a structural reason, not relying on zero locations, why the Lorentzian smoothing of $M_\xi$ against $1-\cos xy$ stays positive in the strip $\frac12<\alpha\le1$. (Positivity for $\alpha\ge1$ is Proposition~\ref{bp:prop:valpha}.)
\end{problem}

\begin{problem}\label{cg:prob:multiplier}
Find a one parameter family of even multipliers $m_\delta(u)\ge1$, $m_\delta\to1$, for which the walk interpretation of Proposition~\ref{cg:prop:spitzer} survives (for instance $m_\delta$ arising from subordinating the stable walk) and for which real zeros of $\int_\R m_\delta\Phi\,e^{iut}du$ can be proved for $\delta$ below the known threshold $t=0.2$ of the flow.
\end{problem}

\begin{problem}\label{cg:prob:cosh}
The symmetric damping $m_\delta(u)=e^{-\delta\cosh2u}$ is even and smooth, so every constraint of Corollary~\ref{cg:cor:modular} holds for $m_\delta\Phi$, and its improved decay is immediate. Decide whether its transforms lie in $\CGLP$. It is not the same as weighting $\theta$ by the Jacobi invariant weight $e^{-\delta(x+1/x)/2}$: with $L$ the operator taking $\theta$ to $\Phi$, $L(W\theta)\neq WL\theta$ because the product rule leaves terms in $W'$ and $W''$, so $m_\delta\Phi$ is the right object. For $A,B>0$ one has $\int_\R e^{(a+it)u-Ae^{2u}-Be^{-2u}}\,du=(B/A)^{(a+it)/4}K_{(a+it)/2}(2\sqrt{AB})$ (substitute $x=e^{2u}$ in the integral defining $K$), and each term of $m_\delta\Phi$ has $A=\pi n^2+\delta/2$, $B=\delta/2$ and $a=\frac92$ or $\frac52$, which gives Bessel $K$ sums of Chowla Selberg type (compare Chapter~\ref{chap:ki}); this identity gives no conclusion about zeros. A height of order $\exp(c/\delta)$ at which nonreal zeros appear is only a conjecture suggested by the multiplier $e^{-\delta\cosh2u}\le1$, with no theorem or numerical certification, and the low lying numerics are not decisive.
\end{problem}

Proposition~\ref{bp:prop:valpha} settles the source's first open problem, compensated positivity of $\xi$ for $1<\alpha<7.789$.
